\documentclass[10pt,a4paper]{amsart}
\usepackage[margin=19mm]{geometry}
\usepackage{amsmath,amssymb,mathtools}
\usepackage[scr=rsfs,cal=euler]{mathalfa}
\usepackage{xfrac}
\usepackage[unicode,hidelinks,bookmarks=false]{hyperref}
\newcommand{\ie}{i.e.\ }
\newcommand{\cf}{cf.\ }
\newcommand{\resp}{respectively\ }
\newcommand{\Subsection}{\subsection}
\providecommand{\nobreakdash}{\nobreak\hbox{-}}
\setlength{\emergencystretch}{2em}
\multlinegap0pt
\usepackage{amssymb}
\usepackage{mathtools}
\usepackage{float}
\usepackage{xcolor}
\usepackage[shortlabels]{enumitem}
\newtheorem{theorem}{Theorem}
\title{Persistence and long-time breakdown of most probable paths under time-dependent fractional noise with applications to KAM tori}
\author{Yanbin Zhu, Xiaomeng Jiang, Yong Li}
\date{Source: arXiv:2609.11645v1 (CC BY 4.0)}
\begin{document}
\maketitle

\noindent\emph{Source and licence.} Persistence and long-time breakdown of most probable paths under time-dependent fractional noise with applications to KAM tori. Version arXiv:2609.11645v1 (CC BY 4.0), distributed by arXiv under the Creative Commons Attribution 4.0 International licence, http://creativecommons.org/licenses/by/4.0/, as recorded in the arXiv licence metadata for this exact version and shown on the abstract page https://arxiv.org/abs/2609.11645v1. Full licence notice: \texttt{licenses/arxiv-2609.11645-CC-BY-4.0.txt}.\par\smallskip

\noindent\textit{Editorial scope.} Counted unit: the article's theorem thm:H-greater-half-collapse (source0.tex lines 3085--3104). The article's complete original proof of it is delivered with it (source0.tex lines 3106--3319, one \texttt{proof} environment of 214 lines). No mathematics is written, repaired or re-derived: the statement and the proof are the article's own lines, in the article's own notation, and no retained line is substituted or re-flowed.  Retained uncounted as the source-local context the proof needs: lines 1773--1780; lines 1963--1968; lines 2913--2930.  Nothing was omitted from any retained span, so the package carries no omission record.  No retained line contains a citation, so the wrapper carries no bibliography and no bibliography entry.  No retained line contains a figure environment, an image-inclusion command or drawing code, so no image is delivered and the package declares no asset dependency.  Editorial formatting only, in addition: an AMS \texttt{amsart} preamble that restates the article's own declarations and macros used by the retained lines, namely the environments \texttt{theorem} on separate counters, together with the standard packages those lines need.  The retained environments print in this document as Theorem 1 in order of appearance, whereas the article prints the same items with the numbering of its own full document.  The complete native source of the article as distributed in the arXiv e-print for this exact version is bundled unchanged as \texttt{sources/210039/source0.tex}.
\begin{equation}
\label{nnoise}
    \bar X_t'
    =
    b_t(\bar X_t),
    \qquad
    \bar X_0=x_0,
\end{equation}

\begin{equation}
\label{cond1}
    \nabla(\nabla\!\cdot b_t)(\bar X_t)=0,
    \qquad
    0\leq t\leq T.
\end{equation}

\begin{equation}
\label{eq:scaling-weighted-fractional-derivative}
\begin{aligned}
&\int_0^{NT}
\left|
    t^\alpha D_{0^+}^\alpha
    \bigl(t^{-\alpha}f(t)\bigr)
\right|^2\,dt
\\
&\qquad=
N^{1-2\alpha}
\int_0^T
\left|
    \tau^\alpha D_{0^+}^\alpha
    \bigl(\tau^{-\alpha}f_N(\tau)\bigr)
\right|^2\,d\tau,
\end{aligned}
\end{equation}

\begin{theorem}[Loss of local minimality on long time intervals for $H>1/2$]
\label{thm:H-greater-half-collapse}
Let $H>1/2$, and let $\bar X$ be a transversally hyperbolic
$T$-periodic solution of \eqref{nnoise} in $\mathbb R^n$, $n\geq2$.
Suppose that \eqref{cond1} holds and that the effective potential has
a positive direction:
\begin{equation*}
    \mathcal V_{\mathrm{eff}}([\omega^*])>0
\end{equation*}
for some
$[\omega^*]\in\widetilde{\mathcal W}_{\mathrm{CR}}$. Then, for every
$\epsilon>0$, there exists
$N_{\mathrm{crit}}(\epsilon)\in\mathbb N^+$ such that, whenever
$N>N_{\mathrm{crit}}(\epsilon)$, the periodic extension of $\bar X$ to
$[0,NT]$ is not a local minimizer of $J_{\epsilon,N}$ in
\begin{equation*}
    x_0+
    \mathcal H_H^\sigma\bigl([0,NT];\mathbb R^n\bigr).
\end{equation*}
\end{theorem}

\begin{proof}
Set $\alpha=H-\frac12\in(0,\frac12)$. For an admissible perturbation
$\eta$, write
\begin{equation}
\label{eq:IHN-def-H-greater-half}
    \mathcal I_H^N(\eta)
    :=
    \int_0^{NT}
    \left|
        s^\alpha D_{0^+}^{\alpha}
        \left[
            s^{-\alpha}\sigma_s^{-1}
            \bigl(
                \eta_s'
                -
                \nabla b(\bar X_s)\eta_s
            \bigr)
        \right]
    \right|^2\,ds
\end{equation}
and
\begin{equation}
\label{eq:VN-def-H-greater-half}
    \mathcal V^N(\eta)
    :=
    -
    \int_0^{NT}
    \eta_s^T
    \nabla^2(\nabla\!\cdot b)(\bar X_s)
    \eta_s\,ds.
\end{equation}
Then
\begin{equation*}
    \delta^2J_{\epsilon,N}(\bar X)[\eta,\eta]
    =
    \frac1\epsilon\mathcal I_H^N(\eta)
    -
    \frac{d_H}{2}\mathcal V^N(\eta).
\end{equation*}

Choose a representative $\omega^*$ of $[\omega^*]$ and set
\begin{equation}
\label{eq:positive-Veff}
    C_V
    :=
    \mathcal V_{\mathrm{eff}}([\omega^*])
    =
    -
    \int_0^T
    (\omega_s^*)^T
    \nabla^2(\nabla\!\cdot b)(\bar X_s)
    \omega_s^*\,ds
    >
    0.
\end{equation}
Since $\omega^*$ is a constant-response mode, there exists a constant
$q^*\in\mathbb R^n$ such that
\begin{equation}
\label{eq:constant-response-selected}
    \sigma_s^{-1}
    \bigl(
        (\omega_s^*)'
        -
        \nabla b(\bar X_s)\omega_s^*
    \bigr)
    =
    q^*.
\end{equation}

Choose a nonzero $\chi\in C^\infty([0,T])$ satisfying
\begin{equation*}
    \chi(0)=\chi(T)=\chi'(0)=\chi'(T)=0
\end{equation*}
and define
\begin{equation}
\label{eq:test-variation-etaN}
    \eta_N(s)
    :=
    \chi\left(\frac{s}{N}\right)\omega_s^*,
    \qquad
    s\in[0,NT].
\end{equation}
Then $\eta_N(0)=\eta_N(NT)=0$, so $\eta_N$ is admissible in both the
fixed-endpoint and free-endpoint perturbation spaces.

We first estimate the potential term. Since
\begin{equation*}
    F(s)
    :=
    (\omega_s^*)^T
    \nabla^2(\nabla\!\cdot b)(\bar X_s)
    \omega_s^*
\end{equation*}
is $T$-periodic and
$\int_0^T F(s)\,ds=-C_V$, periodic averaging gives
\begin{align*}
    \mathcal V^N(\eta_N)
    &=
    -
    \int_0^{NT}
    \chi^2\left(\frac{s}{N}\right)F(s)\,ds
    \\
    &=
    \frac{C_V}{T}
    \left(
        \int_0^T\chi^2(u)\,du
    \right)N
    +
    o(N).
\end{align*}
Hence, for some $C>0$ and all sufficiently large $N$,
\begin{equation}
\label{eq:potential-linear-positive}
    \mathcal V^N(\eta_N)\geq CN.
\end{equation}

For the kinetic term, \eqref{eq:constant-response-selected} and
\eqref{eq:test-variation-etaN} yield
\begin{equation*}
\begin{aligned}
&\sigma_s^{-1}
\bigl(
    \eta_N'(s)-\nabla b(\bar X_s)\eta_N(s)
\bigr)
\\
&\qquad=
\chi\left(\frac{s}{N}\right)q^*
+
\frac1N
\chi'\left(\frac{s}{N}\right)
\sigma_s^{-1}\omega_s^*.
\end{aligned}
\end{equation*}
The scaling estimate
\eqref{eq:scaling-weighted-fractional-derivative} gives
\begin{equation*}
\begin{aligned}
&\int_0^{NT}
\left|
s^\alpha D_{0^+}^{\alpha}s^{-\alpha}
\left[
    \chi\left(\frac{s}{N}\right)q^*
\right]
\right|^2\,ds
\leq
CN^{1-2\alpha}.
\end{aligned}
\end{equation*}

Set $P_*(s):=\sigma_s^{-1}\omega_s^*$. For the remaining term, scaling
gives
\begin{align*}
&\frac1{N^2}
\int_0^{NT}
\left|
s^\alpha D_{0^+}^{\alpha}s^{-\alpha}
\left[
    \chi'\left(\frac{s}{N}\right)P_*(s)
\right]
\right|^2\,ds
\\
&\qquad=
N^{-1-2\alpha}
\int_0^T
\tau^{2\alpha}
\left|
D_{0^+}^{\alpha}
\left[
    \tau^{-\alpha}\chi'(\tau)P_*(N\tau)
\right]
\right|^2\,d\tau.
\end{align*}
Because $\chi'(0)=0$, $P_*$ is smooth and periodic, and
$\alpha<\frac12$, the fractional integral estimate
$D_{0^+}^{\alpha}=I_{0^+}^{1-\alpha}\frac{d}{d\tau}$ gives
\begin{equation*}
    \int_0^T
    \tau^{2\alpha}
    \left|
    D_{0^+}^{\alpha}
    \left[
        \tau^{-\alpha}\chi'(\tau)P_*(N\tau)
    \right]
    \right|^2\,d\tau
    \leq
    CN^2.
\end{equation*}
Consequently,
\begin{equation}
\label{eq:kinetic-upper-H-greater-half}
    \mathcal I_H^N(\eta_N)
    \leq
    CN^{1-2\alpha}.
\end{equation}

Combining \eqref{eq:potential-linear-positive} and
\eqref{eq:kinetic-upper-H-greater-half}, we obtain
\begin{equation}
\label{eq:second-var-final-estimate}
    \delta^2J_{\epsilon,N}(\bar X)[\eta_N,\eta_N]
    \leq
    \frac{C}{\epsilon}N^{1-2\alpha}
    -
    CN.
\end{equation}
Since $\alpha>0$, the first term grows sublinearly, whereas the second
is negative and grows linearly. Thus
\begin{equation*}
    \delta^2J_{\epsilon,N}(\bar X)[\eta_N,\eta_N]<0
\end{equation*}
for all sufficiently large $N$. Therefore, $\bar X$ is not a local
minimizer of $J_{\epsilon,N}$ in
$x_0+\mathcal H_H^\sigma([0,NT];\mathbb R^n)$.
\end{proof}

\end{document}
